% =====================================================================
%  五、模型建立与求解（核心章节）
%  每个子问题统一采用「建模思路 → 模型建立 → 算法/求解 → 结果分析」结构。
% =====================================================================
\section{模型建立与求解}

% ---------------------------------------------------------------------
\subsection{问题一：示向度交会定位模型}
\subsubsection{模型建立}
设第 $i$ 个检测点为 $S_i=(x_i,y_i)$，测得的示向度为 $\theta_i$（$x$ 轴正向逆时针，
单位度），误差界为 $\varepsilon=1^\circ$。由于误差有界，无论真实误差取何值，干扰源都
落在以 $S_i$ 为顶点、关于示向度方向对称张开 $2\varepsilon$ 的角楔之内：
\begin{equation}
\mathcal{W}_i=\bigl\{X:\ \bigl|\beta(S_i,X)-\theta_i\bigr|\le\varepsilon\bigr\},
\label{eq:wedge}
\end{equation}
其中 $\beta(S,X)$ 为自 $S$ 指向 $X$ 的方位角。$m$ 次测向给出的定位区域即各角楔之交。角楔是凸集，也是两个半平面的交；各角楔
对干扰源的约束同时成立，故定位区域为 $\mathcal{R}=\bigcap_{i=1}^{m}\mathcal{W}_i$，
它是 $2m$ 个半平面之交，必为凸多边形，其顶点数不超过 $2m$，且只能落在 $2m$ 条边界
直线的交点上。记该区域的顶点集为 $V$。

凸多边形上任意两点的最大距离必由两个顶点达到，故求直径不必在连续区域上作最优化，
只需在顶点集 $V$ 上比较点对距离：
\begin{equation}
D=\max_{P,Q\in\mathcal{R}}\|P-Q\|=\max_{P,Q\in V}\|P-Q\| .
\label{eq:diameter}
\end{equation}
设使式~\eqref{eq:diameter} 取到最大值的两个顶点为 $A,B$，则直径圆圆心为
$M=(A+B)/2$、半径为 $r=D/2$。

若各示向度方向过于接近，$\mathcal{R}$ 沿公共方向无限延伸，式~\eqref{eq:diameter} 无
定义，故须先判定有界性。由于无界是无穷远处的性质，判据只涉及方向：$\mathcal{R}$ 无界
当且仅当各方位角区间 $\Theta_i=[\theta_i-\varepsilon,\theta_i+\varepsilon]$ 存在公共
交集，当 $m=2$ 时即两次测向的方向夹角不超过 $2\varepsilon$。该判据只需 $O(m)$ 次角度
比较，且不依赖检测点坐标。

覆盖与否由圆盘的凸性直接判定：$\mathcal{R}\subseteq$ 该直径圆当且仅当其全部顶点含于
圆内，即
\begin{equation}
\max_{P\in V}\|P-M\|\le D/2 .
\label{eq:cover}
\end{equation}
条件不满足时，必有顶点 $P$ 使 $\angle APB<90^\circ$ 而落在圆外。注意该圆并把圆心固定
在 $M$ 上，与圆心可自由选取的最小包围圆不同，故式~\eqref{eq:cover} 需逐例检验而不能
直接断言成立。

\subsubsection{模型求解算法}
算法输入为 $m$ 组“检测点坐标与示向度”及误差界 $\varepsilon$，输出为顶点集 $V$、直径
$D$、直径端点 $A,B$、圆心 $M$ 及覆盖判定，流程如图~\ref{fig:p1-flow}，分四步。

第一步有界性预判：对 $m$ 个方位角区间 $[\theta_i-\varepsilon,\theta_i+\varepsilon]$
求交（跨 $0^\circ$ 按圆周区间处理），交集非空即区域无界，直接报“不可用”。该步只作角度
比较、不需坐标，置于最前可避免在发散区域上继续计算。

第二步区域构造：依次用式~\eqref{eq:wedge} 的 $2m$ 条边界所对应的半平面裁剪一个大矩形，
每裁剪一次顶点数至多增加 $1$，故中间结果不超过 $2m$ 个顶点，裁剪结束即得 $\mathcal{R}$。

第三步求直径：取凸包以剔除内部交点与共线中间点，再按式~\eqref{eq:diameter} 枚举顶点对
取距离最大者作为 $D$，对应两点为 $A,B$，圆心 $M=(A+B)/2$。枚举量至多 $O(m^2)$，本题
$m\le20$，亦可改用旋转卡壳在 $O(n)$ 内求解。

第四步覆盖判定：按式~\eqref{eq:cover} 计算 $\max_{P\in V}\|P-M\|$ 并与 $D/2$ 比较，
容差随坐标量级取相对量（本文取 $\tau=10^{-9}\max(1,\max_i\|S_i\|_\infty)$），否则相切
构型会被舍入误差判为不覆盖；若判定不覆盖，同时输出使该值最大的顶点作为越界位置。

\begin{figure}[!htbp]
    \centering
    \includegraphics[width=.58\textwidth]{p1-algorithm-flow}
    \caption{问题一求解算法流程：无界构型在第 1 步即被拦截，有界构型依次经过区域构造、
    求直径与覆盖判定}\label{fig:p1-flow}
\end{figure}

\subsubsection{求解结果与分析}
取五组自造良态算例（检测点分布在干扰源四周，方位张角 $180^\circ\sim250^\circ$、最小
方位间隔不小于 $25^\circ$、与干扰源距离 $160\sim1300$ m），结果见表
~\ref{tab:p1-result}；五组算例的定位区域直径均在 $10\sim26$ m 之间，即两次以上交会已
能把候选范围从测向距离的百米量级收缩到数十米以内。表中同时给出覆盖判定，其两种结果
的对照见图~\ref{fig:p1-coverage}。

\begin{table}[!htbp]
    \caption{问题一算例的计算结果}\label{tab:p1-result}
    \centering
    \begin{tabular}{cccccccc}
        \toprule
        算例 & $m$ & 顶点数 & $D$ (m) & $D/2$ (m) & $\max\|P-M\|$ (m) & 比值 & 覆盖 \\
        \midrule
        $A_1$ & 6 & 6 & 17.64 & 8.82 & 8.82 & 1.0000 & 是 \\
        $A_2$ & 5 & 5 & 18.61 & 9.31 & 10.29 & 1.1055 & 否 \\
        $A_3$ & 4 & 5 & 25.64 & 12.82 & 13.75 & 1.0722 & 否 \\
        $A_4$ & 5 & 4 & 21.12 & 10.56 & 10.56 & 1.0000 & 是 \\
        $A_5$ & 6 & 5 & 10.59 & 5.30 & 5.30 & 1.0000 & 是 \\
        \bottomrule
    \end{tabular}
\end{table}

可见五组算例中有两组（$A_2$、$A_3$）的直径圆无法覆盖定位区域，越界顶点到 $M$ 的距离
分别比 $D/2$ 大 $0.98$ m 与 $0.93$ m，超出 $10.55\%$ 与 $7.22\%$；其原因在于这两组
算例的区域是细长的斜四边形，直径为其一条对角线，而越界顶点“斜看”直径的张角小于
$90^\circ$（$A_2$ 中该顶点处的张角为 $84.2^\circ$），因而落在圆外。

\begin{figure}[!htbp]
    \centering
    \includegraphics[width=.88\textwidth]{p1-coverage-contrast}
    \caption{直径圆覆盖判定的两种结果：(a) 覆盖成立（算例 $A_1$，比值 $1.0000$）；
    (b) 覆盖不成立（算例 $A_2$，比值 $1.1055$），橙色虚线标出落在圆外的顶点 $P$ 及其
    $|PM|=10.29\ \mathrm{m}>D/2=9.31\ \mathrm{m}$}\label{fig:p1-coverage}
\end{figure}

为说明这种失败是否偶然，本文以固定随机种子作两组扫描：在目标区域内随机生成 $20000$
组两测点构型（只保留有界者），覆盖失败率为 $1.16\%$、最坏比值 $1.0336$；随机生成
$4000$ 组 $3\sim6$ 个检测点的良态构型，其中 $3472$ 个有界构型有 $113$ 个失败、失败率
$3.25\%$、最坏比值 $1.0304$，即最坏时直径圆需再放大约 $3\%$ 才能覆盖区域。

综上，以定位区域直径为直径的圆不一定能覆盖该区域，须按式
~\eqref{eq:cover} 逐例判定；覆盖成立时该圆恰为区域的最小包围圆，不成立时其所需半径
比 $D/2$ 略大（算例 $A_2$、$A_3$ 分别大约 $1.005$ 倍与 $1.001$ 倍），最小包围圆与
该圆的差别由越界顶点给出。

% ---------------------------------------------------------------------
\subsection{问题二：第二检测点的选择模型}
\subsubsection{模型建立}
\textrm{【待填充：给出该问的模型与求解方法。】}

\subsubsection{求解与结果}
\textrm{【待填充：给出该问的求解结果与分析。】}

% ---------------------------------------------------------------------
\subsection{问题三：多源定位与清除策略}
\subsubsection{模型建立}
\textrm{【待填充：给出该问的模型与求解方法。】}

\subsubsection{求解与结果}
\textrm{【待填充：给出该问的求解结果与实测指标。】}

% ---------------------------------------------------------------------
\subsection{问题四：含定向源时的定位与清除模型}
\subsubsection{模型建立}
\textrm{【待填充：给出该问的模型与求解方法。】}

\subsubsection{求解与结果}
\textrm{【待填充：给出该问的求解结果与分析。】}
